\scriptsize
\thispagestyle{empty}
\def\subs{\colorea{blue}{A,T,D}}
\begin{frame}

    \begin{align*}
        G_{\subs} & =  \left(S,NT, T,P_\subs\right)        \\
        NT        & = \left\{
                  & C_{11}, C_{12},\ldots, C_{1T}, \ldots, \\
                  & C_{21}, C_{22},\ldots, C_{1T}, \ldots, \\
                  & C_{31}, C_{32},\ldots, C_{3T}, \ldots, \\
        \right\}
    \end{align*}

\end{frame}

% Frame 1: Gramática de Alto Nivel
\begin{frame}\scriptsize
    \begin{tikzpicture}[node distance=0.1cm, every node/.style={align=center}]
        \node[minimum width=6cm] (Grammar)
        {$G_{\subs} =  \left(S,NT, T,P_\subs\right)$};

        % No terminales
        \node[xshift=0.55cm,below=-0.35mm of Grammar, text width=5cm,inner sep=3mm,] (NonTerminalsA)
        {$NT = \left\{
                C_{11}, C_{12},\ldots, C_{1T}, \ldots,
                % C_{21}, C_{22},\ldots, C_{1T}, \ldots,
                % C_{31}, C_{32},\ldots, C_{3T}, \ldots,
                \right. $};
        \node[xshift=0.5cm,below=-0.5cm of NonTerminalsA, text width=5cm,inner sep=3mm,] (NonTerminalsB)
        {$C_{21}, C_{22},\ldots, C_{1T}, \ldots,$};
        \node[xshift=-0.27cm,below=-0.5cm of NonTerminalsB, text width=5cm,inner sep=3mm,] (NonTerminalsC)
        {$C_{31}, C_{32},\ldots, C_{AT},$};
        \node[xshift=-0.35cm,below=-0.5cm of NonTerminalsC, text width=5cm,inner sep=3mm,] (NonTerminalsD)
        {$Z_3,\ldots, Z_A \left.\right\}$};

        % Terminales
        \node[xshift=0.15cm,below=0.5cm of NonTerminalsC, minimum width=6cm] (Terminals)
        {$T = \left\{:, ;, 0, 1, 2,3,,4,5,6,7,8, 9\right\}$};

        % % Encabezado
        \node[xshift=-2.10cm,below =-0.35cm of Terminals, minimum width=6cm, minimum height=1cm] (Axiom)
        {$P_{\subs} =  \left\{ \right.$};
        % % Producciones
        \node[draw, rectangle, below=0.5cm of Terminals, minimum width=6cm] (Productions)
        {
            $
                \begin{aligned}
                    S      & ::= C_{11}, C_{12}: \ldots : C_{1T};         \\
                           & \qquad C_{21}, C_{22}: \ldots : C_{2T};      \\
                           & \qquad C_{11}, C_{12}: \ldots : C_{1T};      \\
                           & \qquad Z_3: C_{31}, C_{32}: \ldots : C_{3T}; \\
                           & \qquad\qquad \vdots                          \\
                           & \qquad Z_A: C_{A1}, C_{A2}: \ldots : C_{AT}; \\
                    Z_3    & ::= 1 \mid 2                                 \\
                           & \qquad\qquad \vdots                          \\
                    Z_A    & ::= 1 \mid 2 \mid \ldots \mid  A-1           \\
                    Z_A    & ::= 1 \mid 2 \mid \ldots \mid  A-1           \\
                    C_{11} & ::= 0 \mid 1 \mid D_1                        \\
                           & \qquad\qquad \vdots                          \\
                    C_{1T} & ::= 0 \mid 1 \mid D_1                        \\
                \end{aligned}
            $
            %     $Z3 \rightarrow 1 \mid 2 \mid 3$ \\
            %     $C1 \rightarrow 0 \mid 1 \mid 2$ \\
            %     $C2 \rightarrow 0 \mid 1 \mid 2 \mid 3 \mid 4$ \\
            %     $C3 \rightarrow 0 \mid 1 \mid 2$
        };

    \end{tikzpicture}
\end{frame}

% \thispagestyle{empty}
% % Frame 2: Gramática de Bajo Nivel
% \begin{frame}{Gramática de Bajo Nivel}
%     \begin{center}
%         \begin{tikzpicture}[node distance=1cm, every node/.style={align=center}]

%             % Encabezado
%             \node[draw, rectangle, minimum width=6cm, minimum height=1cm] (AxiomLL)
%             {\textbf{Axioma:} \\ $S$};

%             % No terminales
%             \node[draw, rectangle, below=0.5cm of AxiomLL, minimum width=6cm] (NonTerminalsLL)
%             {\textbf{No Terminales:} \\ $S, R, Z, R3, R4$};

%             % Terminales
%             \node[draw, rectangle, below=0.5cm of NonTerminalsLL, minimum width=6cm] (TerminalsLL)
%             {\textbf{Terminales:} \\ $:, ;, -, 0, 1, 2, \dots, 4$};

%             % Producciones
%             \node[draw, rectangle, below=0.5cm of TerminalsLL, minimum width=6cm] (ProductionsLL)
%             {\textbf{Producciones:} \\
%                 $S \rightarrow R ; Z ; treasures ; monsters$ \\
%                 $R \rightarrow R3 : R4$ \\
%                 $R3 \rightarrow 0 \mid 1$ \\
%                 $R4 \rightarrow 0 \mid 1 \mid 2$ \\
%                 $Z \rightarrow 1 \mid 2 \mid 3$ \\
%                 $treasures \rightarrow gold \mid silver \mid gems$ \\
%                 $monsters \rightarrow skeleton \mid zombie \mid goblin$
%             };

%         \end{tikzpicture}
%     \end{center}
% \end{frame}

% \begin{frame}
%     % 
%     \begin{tikzpicture}
%         \def\yMax{1.1}
%         \begin{axis}[
%                 width=\textwidth,
%                 axis lines=left, % Coloca los ejes en el medio
%                 xlabel=, % Etiqueta del eje x
%                 ylabel=, % Etiqueta del eje y
%                 y=2.5cm,
%                 % x=0.5cm/3,
%                 xmin=-10, xmax=50.9, % Limites del eje x
%                 ymin=0, ymax=\yMax, % Limites del eje y
%                 xtick={-10,0,10,20,...,50}, % Puntos de referencia en el eje x
%                 ytick={0,0.2,...,\yMax}, % Puntos de referencia en el eje y
%                 domain=-10:50, % Dominio de la función
%                 samples=200, % Resolución de la curva
%                 smooth % Curva suave
%             ]
%             \draw[gray, dashed] (axis cs: 0,0) -- (axis cs: 0,1);
%             \draw[gray, dashed] (axis cs: 10,0) -- (axis cs: 10,1);
%             \draw[gray, dashed] (axis cs: 20,0) -- (axis cs: 20,1);
%             \draw[gray, dashed] (axis cs: 30,0) -- (axis cs: 30,1);
%             \draw[gray, dashed] (axis cs: 40,0) -- (axis cs: 40,1);
%             \draw[draw=blue, line width=1.1mm](axis cs: -10, 1) -- (axis cs: 0, 1) -- (axis cs: 10, 0);
%             \draw[draw=green!90!black, line width=1.1mm](axis cs: 0, 0) -- (axis cs: 10, 1) -- (axis cs: 20,0);% -- (axis cs: 30, 0);
%             \draw[draw=cyan!90!black, line width=1.1mm](axis cs: 10, 0) -- (axis cs: 20, 1) -- (axis cs: 30,0);
%             \draw[draw=orange!90!white, line width=1.1mm](axis cs: 20, 0) -- (axis cs: 30, 1) -- (axis cs: 40,0);
%             \draw[draw=red!90!white, line width=1.1mm](axis cs: 30, 0) -- (axis cs: 40, 1) -- (axis cs: 50,1);
%         \end{axis}
%         \node[above,blue] (a) at (1,3){\textbf{Frío}};
%         \node[above,green!90!white] (b) at (2.9,2.93){\textbf{Templado}};
%         \node[above,cyan!90!black] (c) at (4.5,3){\textbf{Fresco}};
%         \node[above,orange!90!white] (d) at (6,3){\textbf{Cálido}};
%         \node[above,red!90!white] (e) at (8,3){\textbf{Caliente}};
%         \node[gray,below] (f) at (4.5,-0.45){\small Temperatura ($^\circ C$)};
%         \node[gray,left,rotate=90] (g) at (-1.25,2){\small $\mu(x)$};
%     \end{tikzpicture}
% \end{frame}

% \begin{frame}
%     % 
%     \begin{tikzpicture}
%         \def\yMax{1.1}
%         \begin{axis}[
%                 width=\textwidth,
%                 axis lines=middle, % Coloca los ejes en el medio
%                 xlabel=, % Etiqueta del eje x
%                 ylabel=, % Etiqueta del eje y
%                 y=2.5cm,
%                 % x=0.5cm/3,
%                 xmin=36.6, xmax=38.4, % Limites del eje x
%                 ymin=0, ymax=\yMax, % Limites del eje y
%                 xtick={36.8,37.0,37.2,...,38.0}, % Puntos de referencia en el eje x
%                 ytick={0,0.2,...,\yMax}, % Puntos de referencia en el eje y
%                 domain=36.8:38, % Dominio de la función
%                 samples=200, % Resolución de la curva
%                 smooth % Curva suave
%             ]
%             \draw[fill=orange, opacity=0.45,draw=none] (axis cs: 37.21,0) rectangle (axis cs: 39,1);
%             \draw[fill=blue!30!white,draw=none] (axis cs: 36.4,0) rectangle (axis cs: 37.19,1);
%             % \draw[gray, dashed] (axis cs: 10,0) -- (axis cs: 10,1);
%             \draw[gray, dashed, line width=0.6mm] (axis cs: 37.2,0) -- (axis cs: 37.2,1);
%             \node[] (a) at (axis cs: 36.9,0.5){NO FIEBRE};
%             \node[] (B) at (axis cs: 37.45,0.5){FIEBRE};

%             % \draw[gray, dashed] (axis cs: 20,0) -- (axis cs: 20,1);
%             % \draw[gray, dashed] (axis cs: 40,0) -- (axis cs: 40,1);
%             % \draw[draw=blue, line width=1.1mm](axis cs: -10, 1) -- (axis cs: 0, 1) -- (axis cs: 10, 0);
%             % \draw[draw=green!90!black, line width=1.1mm](axis cs: 0, 0) -- (axis cs: 10, 1) -- (axis cs: 20,0);% -- (axis cs: 30, 0);
%             % \draw[draw=cyan!90!black, line width=1.1mm](axis cs: 10, 0) -- (axis cs: 20, 1) -- (axis cs: 30,0);
%             % \draw[draw=orange!90!white, line width=1.1mm](axis cs: 20, 0) -- (axis cs: 30, 1) -- (axis cs: 40,0);
%             % \draw[draw=red!90!white, line width=1.1mm](axis cs: 30, 0) -- (axis cs: 40, 1) -- (axis cs: 50,1);
%         \end{axis}
%         % \node[above,blue] (a) at (1,3){\textbf{Frío}};
%         % \node[above,green!90!white] (b) at (2.9,2.93){\textbf{Templado}};
%         % \node[above,cyan!90!black] (c) at (4.5,3){\textbf{Fresco}};
%         % \node[above,orange!90!white] (d) at (6,3){\textbf{Cálido}};
%         \draw[-Triangle,GithubRed, line width=0.9mm] (3,2.8) -- (3.1,2.0);
%         \node[above,greenLabel] (e) at (3,2.8){Frontera nítida};
%         \node[gray,below] (f) at (4.5,-0.45){\small Temperatura ($^\circ C$)};
%         \draw[fill=white,draw=white] (-1,-0.3) rectangle (0.1,2.8);
%     \end{tikzpicture}
% \end{frame}

% \begin{frame}
%     % 
%     \begin{tikzpicture}
%         \begin{axis}[
%                 width=\textwidth,
%                 axis lines=middle, % Coloca los ejes en el medio
%                 xlabel=, % Etiqueta del eje x
%                 ylabel=, % Etiqueta del eje y
%                 y=2.5cm,
%                 % x=0.5cm/3,
%                 xmin=164.5, xmax=195.5, % Limites del eje x
%                 ymin=0, ymax=1.1, % Limites del eje y
%                 xtick={165,170,...,190.0}, % Puntos de referencia en el eje x
%                 ytick={0,0.2,...,1.1}, % Puntos de referencia en el eje y
%                 domain=165:190, % Dominio de la función
%                 samples=200, % Resolución de la curva
%                 smooth % Curva suave
%             ]
%             \draw[draw=none,opacity=0.45, left color=cyan,right color = blue,opacity=0.5] (axis cs: 170.1,0) rectangle (axis cs: 180,1);

%             \draw[draw=blue, line width=1.1mm](axis cs: 170, 0) -- (axis cs: 180, 1) -- (axis cs: 190, 1);
%             \draw[gray, dashed,very thick] (axis cs: 170,0) -- (axis cs: 170,1);
%             % \draw[gray, dashed] (axis cs: 40,0) -- (axis cs: 40,1);
%             % \draw[draw=green!90!black, line width=1.1mm](axis cs: 0, 0) -- (axis cs: 10, 1) -- (axis cs: 20,0);% -- (axis cs: 30, 0);
%             % \draw[draw=cyan!90!black, line width=1.1mm](axis cs: 10, 0) -- (axis cs: 20, 1) -- (axis cs: 30,0);
%             % \draw[draw=orange!90!white, line width=1.1mm](axis cs: 20, 0) -- (axis cs: 30, 1) -- (axis cs: 40,0);
%             % \draw[draw=red!90!white, line width=1.1mm](axis cs: 30, 0) -- (axis cs: 40, 1) -- (axis cs: 50,1);
%         \end{axis}
%         \draw[Triangle-Triangle,GithubRed, line width=0.4mm] (1.6,2.7) -- (4.37,2.7);
%         \node[above,greenLabel] (e) at (3,2.8){Frontera difusa};
%         \node[gray,below] (f) at (4.5,-0.45){\small Estatura ($cm$)};
%         \node[gray,left,rotate=90] (g) at (-1.0,2.3){\small Pertenencia};
%     \end{tikzpicture}
%     \begin{center}
%         % $\mu_A(x)  =  \begin{cases}
%         %     1 \text{ si, } x \in A,\\
%         %     0 \text{ si, } x \notin A
%         %   \end{cases}
%         %   $
%         $$
%             \begin{aligned}
%                 S & = \left\{\colorea{blue}{x} | x \in U\right\}                        \\
%                 A & = \left\{\colorea{blue}{\left(x, \mu_A(x)\right)} | x \in U\right\}
%             \end{aligned}
%         $$


%         % $$
%         % \begin{aligned}
%         %     \mu_A: \mathcal{F} &\rightarrow{} [0,1]\\
%         %     x &\rightarrow{} \mu_A(x)\\
%         % \end{aligned}
%         % $$
%     \end{center}
%     \begin{tabular}{l}
%         % \colorea{colorA}{\texttt{IF}} $X$ \texttt{IS} $A$ \colorea{colorA}{\texttt{THEN}} $Y$ es $B$\\
%         % \colorea{colorA}{\texttt{IF}} $X_1$ \texttt{IS} $A_1$ \texttt{AND} ... $X_N$ \texttt{IS} $A_N$ \colorea{colorA}{\texttt{THEN}} $Y$ es $B$\\[2em]
%         \colorea{colorA}{\texttt{IF}} \texttt{torMax} es alto y \texttt{angTorMin} es pequeño \colorea{colorA}{\texttt{THEN}} \texttt{lesión}
%     \end{tabular}
% \end{frame}